% Partie : set-theory
% Chapitre : spine
% Section : recursion

\documentclass[../../../include/open-logic-section]{subfiles}

\begin{document}

\olfileid[fr]{sth}{spine}{recursion}
\olsection{Les théorèmes de récursion transfinie}

Il nous faut d'abord confirmer que la \olref[valpha]{defValphas} est une
définition valable. Plus généralement, nous devons démontrer que
toute tentative de définition par récursion transfinie aboutit. Tel
est le but de cette section.

\emph{Attention : ce qui suit est délicat.} L'idée générale est
néanmoins assez simple : la récurrence transfinie et le remplacement
garantissent la légitimité de plusieurs formes de récursion
transfinie.\footnote{Rappelons que les formules et les termes peuvent
contenir des paramètres, sauf indication contraire explicite.}

\begin{defn}
Soient $\tau(x)$ un terme, $f$ une fonction et $\alpha$ un ordinal.
Nous disons que $f$ est une \emph{approximation jusqu'à $\alpha$} de
$\tau$ si et seulement si $\dom{f} = \alpha$ et
$(\forall \beta \in \alpha)f(\beta) =
\tau(\funrestrictionto{f}{\beta})$.
\end{defn}

\begin{lem}[Récursion bornée]\ollabel{transrecursionfun}
Pour tout terme $\tau(x)$ et tout ordinal $\alpha$, il existe une
unique approximation jusqu'à $\alpha$ de $\tau$.
\end{lem}
\begin{proof}
Nous allons montrer que, pour tout $\gamma \leq \alpha$, il existe
une unique approximation jusqu'à $\gamma$.

Établissons d'abord l'unicité. Soient $g$ et $h$ des approximations
respectivement jusqu'à $\gamma$ et jusqu'à $\delta$. Une récurrence
transfinie sur leurs arguments montre que
$g(\beta) = h(\beta)$ pour tout
$\beta \in \dom{g} \cap \dom{h} = \gamma \cap \delta =
\min(\gamma, \delta)$. Ainsi, les approximations sont uniques
lorsqu'elles existent et coïncident sur leur domaine commun.

Pour établir l'existence, appliquons maintenant la récurrence
transfinie simple (\olref[ordinals][opps]{simpletransrecursion}) aux
ordinaux $\delta \leq \alpha$.

La fonction vide est manifestement une approximation jusqu'à
$\emptyset$.

Si $g$ est une approximation jusqu'à $\gamma$, alors
$g \cup \{\tuple{\gamma, \tau(g)}\}$ est une approximation jusqu'à
$\ordsucc{\gamma}$.

Si $\gamma$ est un ordinal limite et si $g_\delta$ est une
approximation jusqu'à $\delta$ pour tout $\delta < \gamma$, posons
$g = \bigcup_{\delta \in \gamma} g_\delta$. C'est une fonction,
puisque les divers $g_\delta$ coïncident partout où leurs domaines se
recoupent. Si $\delta \in \gamma$, alors
$g(\delta) = g_{\ordsucc{\delta}}(\delta) =
\tau(\funrestrictionto{g_{\ordsucc{\delta}}}{\delta}) =
\tau(\funrestrictionto{g}{\delta})$.

Cela achève la démonstration par récurrence transfinie.
\end{proof}

Si nous nous autorisons à définir un \emph{terme} plutôt qu'une
fonction, nous pouvons lever la borne $\alpha$ du résultat précédent.
Dans l'énoncé et la démonstration qui suivent, lorsque $\sigma$ est
un terme, nous posons
$\funrestrictionto{\sigma}{\alpha} =
\Setabs{\tuple{\beta, \sigma(\beta)}}{\beta \in \alpha}$.

\begin{thm}[Récursion générale]\ollabel{transrecursionschema}
Pour tout terme $\tau(x)$, nous pouvons définir explicitement un
terme $\sigma(x)$ tel que
$\sigma(\alpha) = \tau(\funrestrictionto{\sigma}{\alpha})$ pour
tout ordinal~$\alpha$.
\end{thm}

\begin{proof}
Pour chaque $\alpha$, le \olref{transrecursionfun} fournit une unique
approximation jusqu'à $\alpha$, notée $f_\alpha$, de~$\tau$.
Définissons $\sigma(\alpha)$ par
$f_{\ordsucc{\alpha}}(\alpha)$. Alors :
\begin{align*}
	\sigma(\alpha) &=
	f_{\ordsucc{\alpha}}(\alpha) \\&=
	\tau(\funrestrictionto{f_{\ordsucc{\alpha}}}{\alpha}) \\&=
	\tau(\Setabs{\tuple{\beta, f_{\ordsucc{\alpha}}(\beta)}}{\beta \in \alpha}) \\&=
	\tau(\Setabs{\tuple{\beta, f_{\ordsucc{\beta}}(\beta)}}{\beta \in \alpha})\\&=
	\tau(\funrestrictionto{\sigma}{\alpha})
\end{align*}
car $f_{\ordsucc{\beta}}(\beta) =
f_{\ordsucc{\alpha}}(\beta)$ pour tout $\beta < \alpha$, comme
dans le \olref{transrecursionfun}.
\end{proof}
\noindent
Notons que le \olref{transrecursionschema} est un \emph{schéma}. Il ne
faut surtout pas s'attendre à ce que $\sigma$ définisse une
fonction, c'est-à-dire une certaine sorte d'\emph{ensemble} : son
domaine $\dom{\sigma}$ serait alors l'ensemble de tous les
ordinaux, en contradiction avec le paradoxe de Burali-Forti
(\olref[ordinals][basic]{buraliforti}).

Il reste toutefois à montrer que le \olref{transrecursionschema}
justifie bien notre définition des $V_\alpha$. Cela ne saute pas
nécessairement aux yeux, mais deviendra clair avec une dernière
version, simple, de la récursion transfinie.

\begin{thm}[Récursion simple]\ollabel{simplerecursionschema}
Pour tous termes $\tau(x)$ et $\theta(x)$, et pour tout ensemble
$A$, nous pouvons définir explicitement un terme $\sigma(x)$ tel que :
\begin{align*}
	\sigma(\emptyset) &= A\\
	\sigma(\ordsucc{\alpha}) &= \tau(\sigma(\alpha)) &&
		\text{pour tout ordinal }\alpha\\
	\sigma(\alpha) &= \theta(\ran{\funrestrictionto{\sigma}{\alpha}})&&
	\text{si }\alpha\text{ est un ordinal limite}
\end{align*}
\end{thm}

\begin{proof}
Commençons par définir un terme $\xi(x)$ de la manière suivante :
% t \xi(x) = A si x n'est pas une fonction de domaine ordinal ;
% sinon, posons :
\[
	\xi(x) =
	\begin{cases}
			A &\text{si $x$ n'est pas une fonction de domaine ordinal,}\\
			&\hspace{1em}\text{ou si $\dom{x}=\emptyset$ ; sinon :}\\
			\tau(x(\alpha)) & \text{si $\dom{x} = \ordsucc{\alpha}$}\\
			\theta(\ran{x}) & \text{si $\dom{x}$ est un ordinal limite}
	\end{cases}
\]
\footnote{Note de l'édition française : dans l'original, le premier
cas ne vise que les objets qui ne sont pas des fonctions de domaine
ordinal. La fonction vide, dont le domaine est l'ordinal
$\emptyset$, n'y entre donc pas, alors que la preuve utilise
$\xi(\emptyset)=A$. Nous avons ajouté explicitement ce cas de base.}
Par le \olref{transrecursionschema}, il existe un terme $\sigma(x)$ tel
que $\sigma(\alpha) =
\xi(\funrestrictionto{\sigma}{\alpha})$ pour tout ordinal
$\alpha$ ; de plus, $\funrestrictionto{\sigma}{\alpha}$ est une
fonction de domaine $\alpha$. Nous montrons que $\sigma$ possède les
propriétés requises par récurrence transfinie simple
(\olref[ordinals][opps]{simpletransrecursion}).

D'abord, $\sigma(\emptyset) = \xi(\emptyset) = A$.

Ensuite, $\sigma(\ordsucc{\alpha}) =
\xi(\funrestrictionto{\sigma}{\ordsucc{\alpha}}) =
\tau(\funrestrictionto{\sigma}{\ordsucc{\alpha}}(\alpha)) =
\tau(\sigma(\alpha))$.

Enfin, lorsque $\alpha$ est un ordinal limite,
$\sigma(\alpha) = \xi(\funrestrictionto{\sigma}{\alpha}) =
\theta(\ran{\funrestrictionto{\sigma}{\alpha}})$.
\end{proof}
\noindent
Pour justifier la \olref[valpha]{defValphas}, il suffit maintenant de
prendre $A = \emptyset$, $\tau(x) = \Pow{x}$ et
$\theta(x) = \bigcup x$. Nous avons enfin légitimé la définition des
$V_\alpha$ !{}

\end{document}
